% GENERATED FILE. Edit canonical lessons, then run npm run generate:books.
% canonical-source-hash: fnv1a64:d0810ce6bdaddb93
% canonical-lessons: ML-C368-kotati, ML-C368-kesu, ML-C368-vicarana, ML-C368-nyayadhipan, ML-C368-prati

\chapter{In Court}
\label{ch:ml-in-court}
\hlchaptermodality{368}

\begin{chapteropening}
\textbf{By the end of this chapter:} \emph{I can say a court of law, a court case, a trial, a judge and the accused.}

Complete the last lesson of chapter 368: I can say a court of law, a court case, a trial, a judge and the accused.
\end{chapteropening}

\section[\texorpdfstring{\ml{കോടതി} (kōṭati)}{കോടതി (kōṭati)}]{\ml{കോടതി} (kōṭati) — a court of law}

\label{lesson:ML-C368-kotati}



\begin{quote}
Before the new one: say the Malayalam for to kidnap, then the Malayalam for to keep watch on.
\end{quote}

\subsection*{You'll want to know: \ml{കോടതി}}
\textbf{\ml{കോടതി}} — \emph{kōṭati} — \textquotedblleft{}a court of law\textquotedblright{}.

\begin{quote}
\ml{നാളെ} \ml{അച്ഛന്} \ml{കോടതിയിൽ} \ml{പോകണം}.
\end{quote}

\begin{quote}
\emph{nāḷe acchanŭ kōṭatiyil pōkaṇaṁ.}
\end{quote}

\begin{quote}
Tomorrow father has to go to court.
\end{quote}

In court is \ml{കോടതിയിൽ}. A lawyer, \ml{വക്കീൽ}, argues there.

\begin{cousinweb}
The origin of \textbf{\ml{കോടതി}} (\emph{kōṭati}) is not settled. It is the everyday Malayalam word for a court of law, where a \ml{വക്കീൽ}, a lawyer, works.
\end{cousinweb}

\begin{grammarlens}[title={What you've built}]
One more word for telling what happened.
\end{grammarlens}

\subsection*{Your turn}
\begin{itemize}
\raggedright
  \item \emph{Say it:} \emph{kōṭati}
  \item \emph{Say it:} \emph{kōṭati}, once more
  \item \emph{Say it:} \emph{nāḷe acchanŭ kōṭatiyil pōkaṇaṁ.} — Tomorrow father has to go to court.
  \item \emph{Recall:} say the Malayalam for to kidnap, then the Malayalam for to keep watch on, then say \emph{kōṭati} again
\end{itemize}

\begin{tcolorbox}[breakable,colback=teal!4,colframe=teal!35!black,title={Before you move on}]
What does \ml{കോടതി} mean? (\textquotedblleft{}A court of law\textquotedblright{}.) Say it once more.
\end{tcolorbox}
\section[\texorpdfstring{\ml{കേസ്} (kēsŭ)}{കേസ് (kēsŭ)}]{\ml{കേസ്} (kēsŭ) — a court case}

\label{lesson:ML-C368-kesu}



\begin{quote}
Before the new one: say the Malayalam for to keep watch on, then the Malayalam for a court of law.
\end{quote}

\subsection*{You'll want to know: \ml{കേസ്}}
\textbf{\ml{കേസ്}} — \emph{kēsŭ} — \textquotedblleft{}a court case\textquotedblright{}.

\begin{quote}
\ml{ആ} \ml{കേസ്} \ml{പത്ത്} \ml{വർഷം} \ml{നീണ്ടു}.
\end{quote}

\begin{quote}
\emph{ā kēsŭ pattŭ varṣaṁ nīṇṭu.}
\end{quote}

\begin{quote}
That case dragged on for ten years.
\end{quote}

A case is filed: \ml{കേസ്} \ml{കൊടുക്കുക}.

\begin{cousinweb}
\textbf{\ml{കേസ്}} (\emph{kēsŭ}) was borrowed from English \emph{case}, in its legal sense. Like many English nouns, it ends in a short -ŭ in Malayalam.
\end{cousinweb}

\begin{grammarlens}[title={What you've built}]
One more word for telling what happened.
\end{grammarlens}

\subsection*{Your turn}
\begin{itemize}
\raggedright
  \item \emph{Say it:} \emph{kēsŭ}
  \item \emph{Say it:} \emph{kēsŭ}, once more
  \item \emph{Say it:} \emph{ā kēsŭ pattŭ varṣaṁ nīṇṭu.} — That case dragged on for ten years.
  \item \emph{Recall:} say the Malayalam for to keep watch on, then the Malayalam for a court of law, then say \emph{kēsŭ} again
\end{itemize}

\begin{tcolorbox}[breakable,colback=teal!4,colframe=teal!35!black,title={Before you move on}]
What does \ml{കേസ്} mean? (\textquotedblleft{}A court case\textquotedblright{}.) Say it once more.
\end{tcolorbox}
\section[\texorpdfstring{\ml{വിചാരണ} (vicāraṇa)}{വിചാരണ (vicāraṇa)}]{\ml{വിചാരണ} (vicāraṇa) — a trial}

\label{lesson:ML-C368-vicarana}



\begin{quote}
Before the new one: say the Malayalam for a court of law, then the Malayalam for a court case.
\end{quote}

\subsection*{You'll want to know: \ml{വിചാരണ}}
\textbf{\ml{വിചാരണ}} — \emph{vicāraṇa} — \textquotedblleft{}a trial\textquotedblright{}.

\begin{quote}
\ml{കേസിന്റെ} \ml{വിചാരണ} \ml{അടുത്ത} \ml{മാസം} \ml{തുടങ്ങും}.
\end{quote}

\begin{quote}
\emph{kēsinṟe vicāraṇa aṭutta māsaṁ tuṭaṅṅuṁ.}
\end{quote}

\begin{quote}
The trial of the case will begin next month.
\end{quote}

A court holds it: \ml{വിചാരണ} \ml{നടത്തുക}.

\begin{cousinweb}
\textbf{\ml{വിചാരണ}} (\emph{vicāraṇa}) is borrowed from Sanskrit \emph{vicāraṇa}, \textquotedblleft{}consideration, reflexion, discussion\textquotedblright{} (Monier-Williams): a trial is the court's weighing of a case.
\end{cousinweb}

\begin{grammarlens}[title={What you've built}]
One more word for telling what happened.
\end{grammarlens}

\subsection*{Your turn}
\begin{itemize}
\raggedright
  \item \emph{Say it:} \emph{vicāraṇa}
  \item \emph{Say it:} \emph{vicāraṇa}, once more
  \item \emph{Say it:} \emph{kēsinṟe vicāraṇa aṭutta māsaṁ tuṭaṅṅuṁ.} — The trial of the case will begin next month.
  \item \emph{Recall:} say the Malayalam for a court of law, then the Malayalam for a court case, then say \emph{vicāraṇa} again
\end{itemize}

\begin{tcolorbox}[breakable,colback=teal!4,colframe=teal!35!black,title={Before you move on}]
What does \ml{വിചാരണ} mean? (\textquotedblleft{}A trial\textquotedblright{}.) Say it once more.
\end{tcolorbox}
\section[\texorpdfstring{\ml{ന്യായാധിപൻ} (nyāyādhipan)}{ന്യായാധിപൻ (nyāyādhipan)}]{\ml{ന്യായാധിപൻ} (nyāyādhipan) — a judge}

\label{lesson:ML-C368-nyayadhipan}



\begin{quote}
Before the new one: say the Malayalam for a court case, then the Malayalam for a trial.
\end{quote}

\subsection*{You'll want to know: \ml{ന്യായാധിപൻ}}
\textbf{\ml{ന്യായാധിപൻ}} — \emph{nyāyādhipan} — \textquotedblleft{}a judge\textquotedblright{}.

\begin{quote}
\ml{ന്യായാധിപൻ} \ml{സാക്ഷിയോട്} \ml{ചോദ്യങ്ങൾ} \ml{ചോദിച്ചു}.
\end{quote}

\begin{quote}
\emph{nyāyādhipan sākṣiyōṭŭ cōdyaṅṅaḷ cōdiccu.}
\end{quote}

\begin{quote}
The judge asked the witness questions.
\end{quote}

The formal word: \ml{ന്യായം}, justice, + \ml{അധിപൻ}, a chief.

\begin{cousinweb}
\textbf{\ml{ന്യായാധിപൻ}} (\emph{nyāyādhipan}) is built from Sanskrit \emph{nyāya}, \textquotedblleft{}rule, method; a judicial sentence\textquotedblright{} (CDIAL 7614), which gave \ml{ന്യായം} (\emph{nyāyaṁ}), and \emph{adhipa}, \textquotedblleft{}a ruler, a commander\textquotedblright{} (Monier-Williams).
\end{cousinweb}

\begin{grammarlens}[title={What you've built}]
One more word for telling what happened.
\end{grammarlens}

\subsection*{Your turn}
\begin{itemize}
\raggedright
  \item \emph{Say it:} \emph{nyāyādhipan}
  \item \emph{Say it:} \emph{nyāyādhipan}, once more
  \item \emph{Say it:} \emph{nyāyādhipan sākṣiyōṭŭ cōdyaṅṅaḷ cōdiccu.} — The judge asked the witness questions.
  \item \emph{Recall:} say the Malayalam for a court case, then the Malayalam for a trial, then say \emph{nyāyādhipan} again
\end{itemize}

\begin{tcolorbox}[breakable,colback=teal!4,colframe=teal!35!black,title={Before you move on}]
What does \ml{ന്യായാധിപൻ} mean? (\textquotedblleft{}A judge\textquotedblright{}.) Say it once more.
\end{tcolorbox}
\section[\texorpdfstring{\ml{പ്രതി} (prati)}{പ്രതി (prati)}]{\ml{പ്രതി} (prati) — the accused; a defendant}

\label{lesson:ML-C368-prati}



\begin{quote}
Before the new one: say the Malayalam for a trial, then the Malayalam for a judge.
\end{quote}

\subsection*{You'll want to know: \ml{പ്രതി}}
\textbf{\ml{പ്രതി}} — \emph{prati} — \textquotedblleft{}the accused; a defendant\textquotedblright{}.

\begin{quote}
\ml{പ്രതി} \ml{കോടതിയിൽ} \ml{ഒന്നും} \ml{പറഞ്ഞില്ല}.
\end{quote}

\begin{quote}
\emph{prati kōṭatiyil onnuṁ paṟaññilla.}
\end{quote}

\begin{quote}
The accused said nothing in court.
\end{quote}

Also a copy of a book or a paper: \ml{രണ്ട്} \ml{പ്രതി}, two copies.

\begin{cousinweb}
\textbf{\ml{പ്രതി}} (\emph{prati}) is usually traced to Sanskrit \emph{prati}, \textquotedblleft{}towards, against, in opposition to\textquotedblright{} (CDIAL 8540): the accused is the one the case stands against.
\end{cousinweb}

\begin{grammarlens}[title={What you've built}]
One more word for telling what happened.
\end{grammarlens}

\subsection*{Your turn}
\begin{itemize}
\raggedright
  \item \emph{Say it:} \emph{prati}
  \item \emph{Say it:} \emph{prati}, once more
  \item \emph{Say it:} \emph{prati kōṭatiyil onnuṁ paṟaññilla.} — The accused said nothing in court.
  \item \emph{Recall:} say the Malayalam for a trial, then the Malayalam for a judge, then say \emph{prati} again
\end{itemize}

\begin{tcolorbox}[breakable,colback=teal!4,colframe=teal!35!black,title={Before you move on}]
What does \ml{പ്രതി} mean? (\textquotedblleft{}The accused; a defendant\textquotedblright{}.) Say it once more.
\end{tcolorbox}
